% ---
% title: The absolute value of a uniform random variable
% tags: function-of-rv, uniform
% ---

\begin{problem}
	Let $X$ be uniformly distributed on $[-1,1]$, that is, its density is
	$f_X(x)=\frac12$ for $-1\le x\le1$ and $f_X(x)=0$ otherwise. Let $Y=|X|$.
	(a) Find the distribution function $F_Y(y)=P(Y\le y)$ of $Y$.
	(b) Find the density $f_Y(y)$ of $Y$. Which known distribution is it?
	(c) Find $P(Y>0.5)$.
\end{problem}

\begin{solution}
	\idea To find the distribution of a function of a random variable, first find
	its distribution function: write the event $\{Y\le y\}$ as an event about $X$,
	whose density we know, and compute its probability as an area. Then
	differentiate to get the density.

	\step{(a)} Since $-1\le X\le1$, the absolute value satisfies $0\le Y\le1$.
	So $F_Y(y)=0$ for $y<0$ (the event $Y\le y$ is impossible: an absolute value
	is never negative), and $F_Y(y)=1$ for $y\ge1$ (the event is certain).

	Now take $0\le y\le1$. The inequality $|X|\le y$ means that $X$ is at most
	$y$ away from 0, on either side:
	\[ \{Y\le y\}=\{|X|\le y\}=\{-y\le X\le y\}. \]
	Note that \emph{both} signs count: $X=-0.3$ and $X=0.3$ both give $Y=0.3$.
	The probability is the area under the density of $X$ between $-y$ and $y$:
	\[ F_Y(y)=P(-y\le X\le y)=\int_{-y}^{y}\frac12\,dx
		=\Big[\frac x2\Big]_{-y}^{y}=\frac y2-\frac{-y}{2}=\frac{2y}{2}=y. \]
	Geometrically this is a rectangle of width $2y$ and height $\frac12$, so its
	area is $2y\cdot\frac12=y$. Hence
	\[ F_Y(y)=\begin{cases}
			0, & y<0,         \\
			y, & 0\le y\le1, \\
			1, & y>1.
		\end{cases} \]
	Check: $F_Y(0)=0$ and $F_Y(1)=1$, so $F_Y$ is continuous and increases from
	0 to 1, as it should.

	\step{(b)} The density is the derivative of the distribution function. On
	$0<y<1$ we have $F_Y(y)=y$, so $f_Y(y)=F_Y'(y)=1$. Outside $[0,1]$ the
	distribution function is constant, so the density is 0. (The value at the
	single points $y=0$ and $y=1$ does not matter: a single point has
	probability 0.) So
	\[ f_Y(y)=\begin{cases}
			1, & 0\le y\le1,        \\
			0, & \text{otherwise}.
		\end{cases} \]
	This is exactly the density of the uniform distribution on $[0,1]$: $Y$ is
	\emph{uniform on $[0,1]$}. Check the total area: a rectangle of width 1 and
	height 1 has area 1.

	The reason is symmetry. Taking the absolute value folds the negative half
	of $[-1,1]$ onto the positive half. The density $\frac12$ on $[-1,0]$ lands
	on top of the density $\frac12$ on $[0,1]$, and together they give
	$\frac12+\frac12=1$ on $[0,1]$.

	\step{(c)} The event $Y>0.5$ is the complement of $Y\le0.5$, so by (a)
	\[ P(Y>0.5)=1-F_Y(0.5)=1-0.5=0.5. \]
	Check with $X$ directly: $|X|>0.5$ happens when $X>0.5$ or $X<-0.5$. These
	two events are mutually exclusive, so their probabilities add. Each is a
	rectangle of width $0.5$ and height $\frac12$ under the density of $X$
	(the shaded areas in the picture):
	\[ P(|X|>0.5)=P(X<-0.5)+P(X>0.5)=0.5\cdot\tfrac12+0.5\cdot\tfrac12
		=0.25+0.25=0.5. \]
	\begin{center}
		\begin{tikzpicture}[x=2.4cm,y=2.4cm,>=stealth]
			\fill[teal!30] (-1,0) rectangle (-0.5,0.5);
			\fill[teal!30] (0.5,0) rectangle (1,0.5);
			\draw[->] (-1.35,0) -- (1.4,0) node[right] {\small $x$};
			\draw[->] (0,0) -- (0,0.85) node[above] {\small $f_X(x)$};
			\draw[navy,very thick] (-1.35,0) -- (-1,0);
			\draw[navy,very thick] (-1,0.5) -- (1,0.5);
			\draw[navy,very thick] (1,0) -- (1.35,0);
			\draw[navy,thick,densely dotted] (-1,0) -- (-1,0.5);
			\draw[navy,thick,densely dotted] (1,0) -- (1,0.5);
			\foreach \x in {-1,-0.5,0.5,1}
			\draw (\x,0.03) -- (\x,-0.03) node[below] {\scriptsize $\x$};
			\node[below left] at (0,0) {\scriptsize $0$};
			\node[above left] at (0,0.5) {\scriptsize $\tfrac12$};
			\node at (-0.75,0.25) {\scriptsize $0.25$};
			\node at (0.75,0.25) {\scriptsize $0.25$};
		\end{tikzpicture}
	\end{center}

	\begin{answer}
		(a) $F_Y(y)=0$ for $y<0$, $F_Y(y)=y$ for $0\le y\le1$, $F_Y(y)=1$ for $y>1$
		\quad (b) $f_Y(y)=1$ for $0\le y\le1$, and 0 otherwise: $Y$ is uniform on $[0,1]$
		\quad (c) $P(Y>0.5)=0.5$
	\end{answer}
	\begin{remark}
		\textbf{Common mistake:} forgetting the negative values of $X$, that
		is, writing $P(|X|\le y)=P(0\le X\le y)=\frac y2$. This gives
		$F_Y(1)=\frac12$, which is impossible for a distribution function of a
		variable with $Y\le1$. The event $|X|\le y$ is the whole interval
		$[-y,y]$, which has two halves. In general, for a function that is not
		monotone (like $|x|$ or $x^2$ on an interval around 0), always write out
		the full set of $x$ values that give $Y\le y$.
	\end{remark}
\end{solution}
